\section{模型评价与推广改进}
\subsection{模型优点}
\textbf{守恒关系明确，四问具有连续结构。}模型从圆柱薄壳收支出发，分别说明储能项、储水项和界面通量的来源。第一问建立基本框架，第二问引入物性反馈，第三问加入全域阈值事件，第四问加入材料运动和几何变化，能逐项对应题目增加的要求。

\textbf{离散方式与物理结构一致。}共享界面通量使相邻单元的内部交换严格抵消；中心用有限控制体处理，避免直接计算奇异形式。非均匀网格把分辨率配置到梯度更大的表面区域，在保留内部连续分布的同时降低均匀粗网格引起的终点偏差。

\textbf{结果具有多角度数值证据。}论文不仅提供求解器输出，还采用独立贝塞尔级数、空间加密、容差收紧和整体收支作核验。事件定位独立于60 s导出间隔，四位小数展示与未舍入判断分开处理，提高了结果报告的一致性。

\textbf{移动域输出含义清楚。}第四问区分材料位置与实际距离，给出坐标变换和表面条件的半径缩放，并对域外位置留空。结果表反映当前真实几何区域，而非沿用初始半径产生没有材料的虚构数值。

\subsection{模型适用边界}
现有模型是一维径向有效近似，未显式描述潜热、湿分迁移焓及端部效应；空气湿度到材料平衡含水率采用有效边界闭合，后期环境和内部收缩速度也包含假设。因而数值精度高不等同于物理预测已通过实验验证，尤其不宜将本文时长直接作为其他品种、尺度或工况的无条件工艺值。

同时，题给经验密度与参考干质量分工使用，避免了在缺乏完整组分和形变数据时强加不相容约束。若将来获得真实湿密度、局部收缩及吸附关系，应统一升级质量与焓模型，不能只替换某一个密度公式。

\subsection{推广与改进方向}
\textbf{完善表面交换。}获得吸附或解吸等温线后，可先将空气湿度转换为材料平衡值$C_e=F(T,\varphi)$；若需显式考虑表面蒸发，则引入潜热通量
\begin{equation}
h(T_a-T_s)=kT_r(R,t)+L_vj_{\rm evap}.
\end{equation}
这里$j_{\rm evap}$必须是真正的相变质量通量，不能把内部每处$C$下降都当作当地蒸发。

\textbf{扩展形变和工艺控制。}可在测得轴向收缩或局部变形后推广为二维轴对称模型；进一步结合能耗、品质损伤指标构建工艺优化问题。当前有限体积、隐式积分和阈值检测框架可继续复用，但新增参数需有实验来源，新增效果需单独计算验证。
